\subsection{Vektoren als Verschiebungspfeile}
\label{subsec:vector_algebra_vectors_as_displacements}

Die Menge aller \emph{Verschiebungspfeile} \EnglishTerm{displacement arrows}
bildet einen Vektorraum: Verschiebungspfeile lassen sich addieren und mit
Skalaren multiplizieren; die Ergebnisse sind wieder Verschiebungspfeile. Wenn
im Folgenden von Vektoren die Rede ist, sind stets diese Verschiebungspfeile
gemeint.

\begin{DefinitionBox}{Vektor als Verschiebungspfeil}{vector_algebra_vector}
  Abhängig vom betrachteten Vektorraum können Vektoren durch unterschiedliche
  mathematische Objekte realisiert werden. Im Vektorraum der Verschiebungen
  stellen wir einen Vektor als Verschiebungspfeil dar. Diese geometrische
  Realisierung ist durch ihre Länge und ihre Richtung eindeutig bestimmt.

  Vektoren kennzeichnen wir in Gleichungen durch einen Pfeil über dem
  mathematischen Symbol, z. B.
  \begin{center}
    \CachedTikzFigure{content/resources/chapters/vector_algebra/tikz/vector_arrow_notation.tex}
  \end{center}
  Ein Verschiebungspfeil besitzt keine feste Position im Raum: Paralleles
  Verschieben ändert den dargestellten Vektor nicht. Zwei Pfeile mit gleicher
  Länge und gleicher Richtung stellen daher denselben Vektor dar.
  \par\smallskip
  \centering
  \CachedTikzFigure{content/resources/chapters/vector_algebra/tikz/vector_translation.tex}
  \par
\end{DefinitionBox}

Um die gegenseitige Ausrichtung von Vektoren beschreiben zu können, haben wir für die
folgenden Spezialfälle eine besondere Notation.
\begin{KeyPointBox}{Beziehungen zwischen Vektoren}{vector_algebra_vector_relations}
  \begin{itemize}[leftmargin=1.4em,itemsep=0.6em,parsep=0pt,topsep=0pt]
    \item
      \begin{tabularx}{\linewidth}[t]{@{}X
          >{\centering\arraybackslash}p{3.8cm}
          >{\centering\arraybackslash}p{3.8cm}@{}}
        gleich \EnglishTerm{equal}: gleicher Betrag und gleiche Richtung
          & $\vec a=\vec b$
          & \CachedTikzFigure{content/resources/chapters/vector_algebra/tikz/equal_vectors_relation.tex}
      \end{tabularx}
    \item
      \begin{tabularx}{\linewidth}[t]{@{}X
          >{\centering\arraybackslash}p{3.8cm}
          >{\centering\arraybackslash}p{3.8cm}@{}}
        parallel \EnglishTerm{parallel}: gleiche Richtung
          & $\vec a\upuparrows\vec b$
          & \CachedTikzFigure{content/resources/chapters/vector_algebra/tikz/parallel_vectors_relation.tex}
      \end{tabularx}
    \item
      \begin{tabularx}{\linewidth}[t]{@{}X
          >{\centering\arraybackslash}p{3.8cm}
          >{\centering\arraybackslash}p{3.8cm}@{}}
        antiparallel \EnglishTerm{antiparallel}: entgegengesetzte Richtung
          & $\vec a\updownarrows\vec b$
          & \CachedTikzFigure{content/resources/chapters/vector_algebra/tikz/antiparallel_vectors_relation.tex}
      \end{tabularx}
    \item
      \begin{tabularx}{\linewidth}[t]{@{}X
          >{\centering\arraybackslash}p{3.8cm}
          >{\centering\arraybackslash}p{3.8cm}@{}}
        \emph{Gegenvektoren} \EnglishTerm{opposite vectors}: gleicher Betrag
        und entgegengesetzte Richtung
          & $\vec b=-\vec a$
          & \CachedTikzFigure{content/resources/chapters/vector_algebra/tikz/opposite_vectors_relation.tex}
      \end{tabularx}
  \end{itemize}
\end{KeyPointBox}
